% !TEX root = ../notes.tex

\documentclass[../notes.tex]{subfiles}

\begin{document}

\section{March 18}
Here we go.

\subsection{Admissible Representations}
For any representation $M$ of $K$, note that there is a canonical map
\[\bigoplus_{\rho\in\op{IrRep}K}\rho\otimes\op{Hom}_K(\rho,M)\to M\]
given by composition, and we showed on the homework that this map is an isomorphism if and only if $M$ is locally finite. This decomposition allows us to define admissibility.
\begin{definition}[admissible]
	Fix a real reductive group $G$ with Lie algebra $\mf g$ and maximal compact subgroup $K$. Then a $(\mf g,K)$-module $M$ is \textit{admissible} if and only if $\dim\op{Hom}_K(\rho,M)$ is finite-dimensional for all irreducible representations $\rho$ of $K$.
\end{definition}
\begin{remark}
	On the homework, we showed that a representation $M$ of a totally disconnected group $G$ is admissible if and only if $\dim\op{Hom}_K(\rho,M)$ is finite-dimensional for all irreducible representations $\rho$ of a choice of compact open subgroup $K$ of $G$.
\end{remark}
We also remark that we can bring down the definition of unitary to $(\mf g,K)$-modules.
\begin{definition}[unitary]
	Fix a real reductive group $G$ with Lie algebra $\mf g$ and maximal compact subgroup $K$. Then a $(\mf g,K)$-module $M$ is \textit{unitary} if and only if $M$ admits a positive-definite Hermitian form $\langle-,-\rangle$ which is invariant under $\mf g$ and $K$.
\end{definition}
\begin{remark}
	Invariance for $K$ is as expected, but invariance for $\mf g$ must be differentiated: by the product rule, it turns out that we are asking for
	\[\langle Xu,v\rangle+\langle u,Xv\rangle=0\]
	for all $X\in\mf g$ and $u,v\in M$.
\end{remark}

\subsection{The Infinitesimal Character}
The center of $U\mf g$ is interesting to us because it acts by scalars on simple modules, and scalars are potentially easier to understand.
\begin{definition}[infinitesimal character]
	Fix a real reductive group $G$ with Lie algebra $\mf g$ and maximal compact subgroup $K$. For a simple $(\mf g,K)$-module $M$, we define the \textit{infinitesimal character} $\chi\colon Z(U\mf g)\to\CC$ by defining
	\[\chi(z)\coloneqq z|_M.\]
	Because $M$ is a $\mf g$-module, $\chi$ is an algebra homomorphism.
\end{definition}
\begin{remark}
	The fact that $z|_M$ is a scalar is some variant of Schur's lemma; one uses the same argument as in \Cref{prop:schur}.
\end{remark}
\begin{remark}
	A choice of character $\chi\colon Z(U\mf g)\to\CC$ is equivalent to giving a maximal ideal of $\op{Sym}(\mf h)^{(W,\cdot)}$, which is equivalent to the data of a (closed) point in $\mf h^*/(W,\cdot)$, which is equivalent to the data of a $(W,\cdot)$-orbit in $\mf h^*$.
\end{remark}
The infinitesimal character turns out to be a very strong invariant.
\begin{proposition}
	Fix a real reductive group $G$ with Lie algebra $\mf g$ and maximal compact subgroup $K$. For any infinitesimal character $\chi\colon Z(U\mf g)\to\CC$, there are only finitely many simple $(\mf g,K)$-modules (up to isomorphism) with this infinitesimal character.
\end{proposition}
\begin{example}
	Let's compute some infinitesimal characters of the simple $(\mf{sl}_2,\op{SO}(2))$-modules $M$. Indeed, let $M=\bigoplus_nM_n$ be the weight decomposition. Then $ef$ should act by some scalar $c_{n-1}$ on $M_n$, and $fe$ should act by some scalar $c_{n+1}$ on $M_n$ as well. Because $[ef,fe]=h$, we see that $c_{n-1}-c_{n+1}=n$ for all $n$. On the other hand, the Casimir $\Omega$ must act by a single eigenvalue $\lambda$ on all the weight spaces $M_n$, so
	\[\lambda=\frac12n^2+c_{n-1}+c_{n+1}\]
	for all $n$. Thus, a choice of $\lambda$ allows us to fully determine the scalars up to a choice of just one scalar; for example, if we want our weights to be bounded below, we could assert that our lowest weight is given at $k$, and then $c_{k-1}=0$, and then we can solve for the remaining weights.
\end{example}
It will turn out that all simple $(\mf g,K)$-modules for $\op{SL}_2(\RR)$ can be globalized to an irreducible representation. The easiest way to show this is by hand. A large source of our representations will come from the principal series.
\begin{definition}[principal series]
	Let $B$ be the upper-triangular Borel subgroup of $\op{SL}_2(\RR)$. For $\varepsilon\in\{0,1\}$ and $s\in\CC$, we define the \textit{principal series} as the induced representation $\op{Ind}_B^G\left(\mathrm{sgn}^\varepsilon\left|\cdot\right|^s\right)$, where $\mathrm{sgn}^\varepsilon\left|\cdot\right|^s$ is viewed as a character of $\RR^\times$, identified with the diagonal torus in $B$.
\end{definition}
\begin{remark}
	Note that $\op{diag}(-1,-1)$ acts by $(-1)^\varepsilon$ on the principal series. Thus, if $\varepsilon=0$, then the corresponding $(\mf g,K)$-module has even weights; conversely, if $\varepsilon=1$, then the corresponding $(\mf g,K)$-module has odd weights. One can also compute the action by Casimir; one finds that $\Omega$ acts by $\frac{1-s^2}2$.
\end{remark}
\begin{remark}
	If $s$ is purely imaginary, then $\mathrm{sgn}^\varepsilon\left|\cdot\right|^s$ is a unitary representation of $B$, so the principal series representation is unitary. It also turns out that the principal series is unitary when $s\in(-1,0)\cup(0,1)$; this complicates the representation theory, but it will not complicate the theory of automorphic forms because they do not appear.
\end{remark}
Most of the remaining $(\mf g,K)$-modules are found inside the discrete series. By way of motivation, recall the following result from the representation theory of compact groups.
\begin{theorem}[Peter--Weyl]
	Fix a compact Lie group $G$. Then the matrix coefficient map
	\[\bigoplus_{V\in\op{IrRep}G}V\boxtimes V^*\to L^2(G)\]
	has dense image.
\end{theorem}

\subsection{The Discrete Series}
We would like a similar decomposition to work for our general Lie groups, so we simply define the discrete series representations to be those which appear.
\begin{definition}[discrete series]
	Fix a Lie group $G$. Then we define a \textit{discrete series representation} of $G$ to be an irreducible representation which appears in $L^2(G)$.
\end{definition}
\begin{remark}
	It is equivalent to ask for $V$ to admit nonzero $v\in V$ and $\ell\in V^*$ for which the matrix coefficient $g\mapsto\ell(gv)$ is in $L^2(G)$.
\end{remark}
Of course, it is now difficult to actually construct the discrete series. For example, even if $V$ is finite-dimen\-sional, we expect $\dim V^*$ total copies of $V$ in $L^2(G)$, so it is going to be hard to find particular copies of $V$.
\begin{example}
	We return to $G=\op{SL}_2(\RR)$. Then the discrete series $D_k^-$ and $D_k^+$ are Verma modules, whose weights are concentrated in some ray. Accordingly, we are on the hunt for a copy of the tensor product of $D_k^-$ with a line of weight $k$ in $L^2(G)$. Let's look for this line in $L^2(G)$. Notably, our line should have weight $k$ for the $\op{SO}(2)$-action, and it should be killed by the weight-lowering operator $f$; these properties generate copies of $D_k^+$. Thus, we are on the hunt for a function $\varphi$ for which
	\[\varphi\left(g\begin{bmatrix}
		\cos\theta & \sin\theta \\ -\sin\theta & \cos\theta
	\end{bmatrix}\right)=e^{ik\theta}\varphi(g),\]
	and with $f=\frac12\begin{bsmallmatrix}
		1 & -i \\ -i & 1
	\end{bsmallmatrix}$, we see that we are asking to vanish under some differential operator. For explicit calculations, it is helpful to identify $G/K$ with the upper-half plane $\mc H$ given by the action of fractional linear transformations (where the stabilizer of $i$ is $K$). Once we do this, we see that vanishing under $f$ is equivalent to vanishing under $\overline\del_z$, so we are hunting for holomorphic functions $\varphi$.
\end{example}
The answer to the construction in the previous example is that $V$ should be the space of holomorphic functions $\varphi\colon\mc H\to\CC$ satisfying the functional equation
\[\left(\begin{bmatrix}
	a & b \\ c & d
\end{bmatrix}f\right)(z)\coloneqq(cz+d)^{-k}\varphi\left(\frac{az+b}{cz+d}\right).\]
In order to have a Hilbert space, we should also try to satisfy the (invariant) condition
\[\int_{\mc H}\left|\varphi(z)\right|^2y^k\cdot\frac{dx\,dy}{y^k}<\infty\]
in order to be able to define a norm. These are called the anti-holomorphic discrete series, for some reason.
\begin{remark}
	It turns out that these discrete series representations exist if and only if $G$ admits a compact maximal torus. (Equivalently, we are asking for $G$ to be an inner form of the compact form.) For example, there are no discrete series representations for $\op{SL}_n(\RR)$ as soon as $n\ge3$. On the other hand, the representations $\op{SU}(p,q)$ do have discrete series. When they exist, there is a classification of the discrete series representations.
\end{remark}

\end{document}